\appendix
\section{Reproducibility and Audit Appendix}
\label{sec:reproducibility-appendix}

This appendix documents the audit checks used in the case study.  It reuses frozen main-workspace artifacts and performs no model training or inference.  It does not recompute any experiment metric, run a new bootstrap or threshold sensitivity, or read Paris formal or protected target content.  The tables separate facts stored in the artifacts from suggestions that may be useful elsewhere.  When a field cannot be established from a frozen artifact, it is reported as \emph{unknown/not claimed} rather than filled in by assumption.

The appendix does not add analyses that were not run.  There was no independent-coder classification exercise, inter-rater agreement statistic, alternative bootstrap block-length analysis, or new Global Fixed uncertainty interval.  The 0.5/1/2\% comparison is arithmetic context only; it does not recompute metrics or change a decision.  The S1--S6 registry is a deterministic single-pipeline classification of one inventory.  Historical project records used a review label for what was actually a separately instantiated automated differential artifact audit, not a blinded external human replication or an inter-rater study.

\subsection{Deterministic S1--S6 decision tree}
\label{sec:appendix-status-tree}

The classifier in \path{experiments/00_registry_reclassification/scripts/rebuild_status_registry.py} applies the following ordered rules to each frozen manifest row.  Order is part of the definition; a later numerical rule cannot overwrite an earlier execution fact.

\begin{enumerate}[label=\textbf{D\arabic*.}]
    \item If \texttt{execution\_status=failed}, assign S5 and require any future attempt to use a new versioned retry identifier.
    \item Otherwise, if \texttt{execution\_status} is not \texttt{complete}, assign S6 and require any future execution to use a new versioned run identifier.
    \item Otherwise, if the row maps to UrbanEV source Table 4, assign S4 because the source paper and repository do not identify the disclosed node uniquely.
    \item Otherwise, if no source comparison cell can be mapped uniquely, assign S4.
    \item Otherwise, collect the disclosed RMSE, MAPE, RAE, and MAE grades and take the worst valid grade.  The frozen bands are exact (relative error no greater than 2\%), near (2--5\%), and trend (greater than 5\%); an empty valid-grade set defaults to trend when a comparison record exists.
    \item If comparison fold coverage is below six, assign S3.  Otherwise, if the worst grade is trend, assign S3.
    \item Otherwise, if the track is \texttt{official} and the worst grade is exact, assign S1.  Every other completed, uniquely mapped, non-trend row is S2.
\end{enumerate}

The rules therefore apply in this order: S5/S6 $\rightarrow$ S4 $\rightarrow$ S1/S2/S3.  A completion mark indicates process completion only; it is not a synonym for S1.  The following examples are the first row of each final status in the authoritative CSV.  The \texttt{classification\_reason} text is an English rendering of the Chinese source field so that the appendix compiles under the paper's Latin-font setup.

\begin{table*}[htbp]
\centering
\caption{One frozen registry example per evidence status.  Full run identifiers and reasons are retained to make the classification path inspectable.}
\label{tab:status-examples}
\scriptsize
\begin{tabular}{>{\raggedright\arraybackslash}p{0.045\linewidth}>{\raggedright\arraybackslash}p{0.36\linewidth}>{\raggedright\arraybackslash}p{0.46\linewidth}}
\toprule
State & Example ID & Reason \\
\midrule
S1 & \path{official-compat-v1-t3-lo-occ-3-1-region-None-42-20-eae54b2d} & Official-compatible protocol is identifiable, all six folds are complete, and the relative error of each of four disclosed metrics is no greater than 2\%. \\
S2 & \path{audited-t3-lo-occ-3-1-region-None-42-20-efa1ea0d} & Six folds and the main protocol align, but an audited correction is present or at least one metric lies in the 2--5\% approximation band. \\
S3 & \path{audited-t3-ar-occ-3-1-region-None-42-20-3b0793b0} & At least one paper-disclosed metric has relative error above 5\%; only trend-level evidence is retained. \\
S4 & \path{audited-t4-lo-occ-3-1-194-None-42-20-42da10db} & Neither the paper nor repository discloses the Table 4 node identity; node 194 is a predeclared reverse-identification candidate, not a unique truth. \\
S5 & \path{audited-t3-lstm-occ-6-3-region-None-42-20-c735895e} & The frozen configuration was executed but failed; the failure artifact is retained and cannot contribute a paper number. \\
S6 & \path{audited-t3-TimesNet-occ-3-3-region-None-42-1-1687eb9d} & The configuration is frozen in the manifest but was not executed or registered after accumulated failures in the same model batch. \\
\bottomrule
\end{tabular}
\end{table*}

\FloatBarrier

\subsection{Model identity is a claim field}
\label{sec:appendix-model-identity}

A model name alone does not identify the exact implementation.  The same family label may refer to an upstream release, a local source snapshot, an existing prediction artifact, a retrained adaptation, or a fixed combination of other artifacts.  Table~\ref{tab:model-identity} reports the identity recorded in the frozen artifacts and lists fields that remain unknown.  This is especially important for Paris: \texttt{Paris\_CAPER\_phase\_only} and \texttt{TimeXer\_local\_audited\_compact\_L168} are development-only local adaptations, not exact transfers of the UrbanEV experts or claims about every upstream implementation.

\subsection{Fail-closed state machine}
\label{sec:appendix-state-machine}
\label{app:fail-closed}

The audit states are:

\begin{center}
\texttt{REGISTERED} $\rightarrow$ \texttt{RUNNING} $\rightarrow$ \texttt{CLOSED\_UNVERIFIED} $\rightarrow$ \texttt{AUDIT\_PASS} $\rightarrow$ \texttt{METRIC\_UNLOCKED} $\rightarrow$ \texttt{GATE\_PASS} or \texttt{GATE\_FAIL}.
\end{center}

Three rules prevent a defective artifact from moving forward.  First, an incomplete matrix, target/order mismatch, invalid control, boundary violation, or provenance failure moves \texttt{CLOSED\_UNVERIFIED} to \texttt{BLOCKED} or \texttt{REPAIR\_REQUIRED}; aggregate metrics remain inaccessible.  Second, a permitted repair creates \texttt{REPAIRED\_UNVERIFIED}, records parent and child hashes and whether predictions changed, and returns to the full audit.  It does not move directly to metric unlock.  If selective repair could use information about which controls failed, the full control family is invalidated and rerun.  Third, \texttt{GATE\_FAIL} is terminal for the current protocol version.  Continuing with a changed model, threshold, loss, or evidence role requires a new version and a still-valid data boundary.

Two incidents show how these transitions worked.  In M7-C, the first audit detected three identity permutations at 69/72 cells while \texttt{metrics\_unlocked=false}.  All 24 shuffled controls were invalidated and rerun before the 72/72 audit passed.  In Paris, 32 prediction jobs closed without evaluation metrics.  Missing-target representation was repaired in 24 artifacts (10,236 entries) without changing predictions, and development metrics were unlocked only after the 32/32 prediction audit passed.  These are integrity events, not accuracy improvements.

\subsection{Common lock schema and invariants}
\label{sec:appendix-lock-schema}

The project used several lock and receipt formats rather than one universal schema.  Table~\ref{tab:lock-schema} gives a common interface for the rules that connect each claim to its evidence.  It does not imply that every historical artifact contains every field.  A missing field must be recorded as \texttt{unknown} or fail the relevant gate; it must not be inferred from a filename.

\begin{table*}[htbp]
\centering
\caption{Normalized lock schema.  Existing examples use different field names; the rightmost column states the rule an implementation must enforce.}
\label{tab:lock-schema}
\scriptsize
\begin{tabular}{>{\raggedright\arraybackslash}p{0.12\linewidth}>{\raggedright\arraybackslash}p{0.21\linewidth}>{\raggedright\arraybackslash}p{0.23\linewidth}>{\raggedright\arraybackslash}p{0.26\linewidth}}
\toprule
Group & Required fields & Existing artifact exemplar & Invariant \\
\midrule
Protocol & \texttt{protocol\_id}, version, status, protocol SHA-256 & router, latency, distillation, and Paris teacher protocol files & Data role, endpoint, reference, gate, and stop rule are immutable before aggregate access. \\
Configuration & \texttt{run\_id}, config fingerprint, model identity, seed, fold, horizon & status registry and run manifests & One identifier maps to one configuration; retries retain parentage rather than overwriting identity. \\
Impl. & code/source hashes; dependency and bundle hashes & Chronos model lock; Paris implementation fingerprints & The implementation that emitted predictions equals the registered implementation or has an explicit non-training amendment receipt. \\
Environment & Python and library versions; accelerator, runtime, precision; optional environment hash & latency hardware/software manifests; Chronos environment block & Scope is explicit.  If no unified environment hash exists, reproducibility is limited rather than implied. \\
Data boundary & source/data hash, role, permitted metadata, fold dates & baseline protocol; Paris boundary lock and source manifest & Development, formal, and protected roles cannot be silently exchanged. \\
Target identity & target-value, index/order, time, spatial-unit, and mask hashes; fold--horizon cell & expert diagnostics, run status, prediction audits & Candidate and reference use identical target positions and semantics for every compared cell. \\
Prediction & prediction SHA-256, shape, finiteness/range, attempt path & per-run receipts and prediction audits & The audited bytes are the bytes aggregated; a repair changing predictions requires new scientific evidence. \\
Metric & metric script/hash, endpoint definition, output hash & metric-lock scripts, aggregate summaries & The metric implementation and weighting rule are fixed; output cannot precede unlock. \\
Unlock & lock time, conditions, authority, receipt SHA-256, matrix/state hash & Paris metrics-unlock receipt; M7-C formal audit & Every condition is conjunctive and machine-checkable; role-specific unlock does not unlock another role. \\
Access & analytical-access counter, operation definition, log/receipt location & Paris protected guard and prediction audit & Counter semantics and time origin are stated; storage migration is recorded separately. \\
Repair & parent hash, child hash, reason, prediction-changed flag, audit return state & shuffled-control amendment and masked-target repair ledger & Repair is traceable and returns to unverified state; non-selective scope is used where observation could bias repair. \\
Decision & reference, point effect, interval, fold/cell/horizon checks, gate vector, authorized next action & router, Chronos, D1, and Paris teacher summaries/decisions & A failed required component yields terminal failure for that version even when another subcheck passes. \\
Claim & claim ID, authoritative artifact/hash, evaluation type, allowed/prohibited wording, scope limit & Table~\ref{tab:claim-trace} & Manuscript wording cannot exceed the evidence type or data role. \\
\bottomrule
\end{tabular}
\end{table*}

\FloatBarrier

\subsection{Operational definition of analytical access}
\label{sec:appendix-access-definition}

The Paris guards distinguish three operations.

\begin{description}
    \item[Storage materialization.] Copying, moving, splitting, preserving, or hashing target-bearing files without using their values.  Creating or checking a path/size/hash manifest also belongs to this category.  These operations are recorded as storage events and do not count as target analysis.
    \item[Approved metadata handling.] Reading only the registered non-target fields (date boundary, station identifier, latitude, longitude, postcode, and area), row counts, file sizes, and hashes needed to establish the split and storage receipt.  State and target values are not used.
    \item[Analytical access.] Any post-lock use of formal or protected target values.  This includes displaying, summarizing, plotting, filtering by, learning from, scoring against, or making a scientific claim from those values.  Manual viewing also counts.  Use in training, validation, model selection, metrics, threshold selection, error analysis, or a paper table counts even when no file is modified.
\end{description}

Under this definition, \texttt{analytical\_access\_count=0} means that no protected label was used after the boundary lock for the listed analytical operations.  It does not mean that target bytes were never materialized: one storage migration created vault shards and moved recoverable sources.  It also does not erase the logged display of two example rows during pre-lock schema discovery.  The counter is supported by project receipts and audits.  It is not enforced by the operating system, does not provide encryption, and cannot prevent deliberate bypass by an authorized user.

These controls are intended to prevent accidental use, or use chosen after earlier results were seen, within the project workflow.  Examples include a script path that reaches a future role, a stale cache, recoverable Git objects, or a metric opened before the registered closure.  The vault path, manifest, hashes, guards, and negative main-project scan reduce these risks.  Operating-system compromise, deliberate bypass, encryption, and immutable access logging are outside scope.  The supported claim is source separation and zero recorded analytical use, not hardened security.

\subsection{Paris fold-synchronous bootstrap, including the partial block}
\label{sec:appendix-paris-bootstrap}

The frozen implementation in \path{experiments/09_distill_v2/scripts/aggregate_paris_teacher.py} used 5,000 replicates and seed 42.  For each development fold, it found target timestamps shared by all four horizons and sorted them.  It required one contiguous hourly sequence and divided that sequence at offsets 0, 24, 48, and so on.  The final slice was kept even when it contained fewer than 24 hours.

Each replicate sampled the fold's original number of blocks with replacement.  The same sampled timestamp sequence was used for every horizon in that fold.  Observed-target masks were then applied within each cell, and cell RMSEs were macro-averaged.  The better CAPER/TimeXer reference was selected again inside the replicate, and teacher gain was computed against that reference.  The reported interval uses the 2.5th and 97.5th percentiles.

August and October each contain 732 common hourly timestamps: 30 full 24-hour blocks and one 12-hour block (31 selectable blocks).  September and November each contain 708 timestamps: 29 full blocks and one 12-hour block (30 selectable blocks).  Every block, including the 12-hour partial block, had the same probability of selection.  Because the partial block contained fewer timestamps, replicate length varied with the number of times it was drawn.  No padding, wraparound, discarding, or extra weighting was applied.  The reported Paris interval applies to this exact rule.  No alternative block length or partial-block sensitivity analysis was run.

\subsection{Claim-to-artifact links}
\label{sec:appendix-claim-lineage}
\label{app:claim-lineage}

Table~\ref{tab:claim-trace} links each headline result to the machine-readable artifact that supports it.  The hash is computed over the named artifact in the main workspace.  A hash identifies file bytes; it does not prove that the method is correct.  The evaluation-type and wording columns state the claim limit.  For example, the Global Fixed row supports an internal point comparison and a separate measured latency result.  It does not gain an unreported bootstrap interval merely because later stages used one.

\begin{table*}[htbp]
\centering
\caption{Claim-to-artifact trace and wording guardrails (part 1 of 2). Hashes are the first 12 hexadecimal characters of SHA-256 over the named frozen artifact.}
\label{tab:claim-trace}
\scriptsize
\begin{tabular}{>{\raggedright\arraybackslash}p{0.04\linewidth}>{\raggedright\arraybackslash}p{0.17\linewidth}>{\raggedright\arraybackslash}p{0.22\linewidth}>{\raggedright\arraybackslash}p{0.11\linewidth}>{\raggedright\arraybackslash}p{0.25\linewidth}}
\toprule
ID & Headline result & Authoritative artifact / SHA-256$_{12}$ & Evaluation type & Allowed wording; prohibited wording \\
\midrule
C1 & Frozen registry: 1,224 configurations; 865 complete artifacts; S1=156. & \path{experiments/00_registry_reclassification/outputs/CONFIG_STATUS_SUMMARY.json}\newline \texttt{833352c136a4} & registry audit & \textit{Allowed:} In this inventory, completion and exact reproduction differ. \newline \textit{Prohibited:} A field-wide reproducibility rate or a model-quality ranking. \\
C2 & Source-table status distributions are uneven under the frozen mapping. & \path{experiments/00_registry_reclassification/outputs/CONFIG_STATUS.csv}\newline \texttt{4b61c2aea510} & registry audit & \textit{Allowed:} Report S1--S6 counts by UrbanEV source table. \newline \textit{Prohibited:} Call an S4 model invalid or intrinsically weak. \\
C3 & Target-dependent CAPER--TimeXer oracle gain is 16.926\%. & \path{experiments/01_expert_diagnostics/outputs/EXPERT_DIAGNOSTIC_SUMMARY.json}\newline \texttt{53fa7f22fca6} & real-GT oracle diagnostic & \textit{Allowed:} An infeasible upper bound indicates local complementarity. \newline \textit{Prohibited:} The oracle is deployable or its gain is recoverable. \\
C4 & State-feature expert-win AUC is 0.6151. & \path{experiments/01_expert_diagnostics/outputs/STATE_FEATURE_DECISION.json}\newline \texttt{9f84ac483e02} & self-supervised proxy & \textit{Allowed:} Relative expert advantage is moderately predictable. \newline \textit{Prohibited:} AUC proves lower forecast RMSE or causal demand knowledge. \\
C5 & Strict Global Fixed RMSE is 0.071367, a 3.178\% internal gain. & \path{experiments/03_fixed_fusions/outputs/FIXED_FUSION_SUMMARY.json}\newline \texttt{0a0162204c83} & internal rolling real GT & \textit{Allowed:} Improves on TimeXer under the audited UrbanEV protocol. \newline \textit{Prohibited:} Blind-test, cross-city, efficient, or deployment-ready. \\
C6 & Direct-loss router gains 0.443\% and fails the frozen 1\% gate. & \path{experiments/05_router_v1/outputs/ROUTER_V1_SUMMARY.json}\newline \texttt{a9dfc2ebe4a6} & internal real GT; gated & \textit{Allowed:} A positive sub-threshold internal signal; current stage stopped. \newline \textit{Prohibited:} Router qualification, robustness, or post-hoc gate revision. \\
C7 & Global Fixed batch-1 P50 is 48.86\% above TimeXer. & \path{experiments/06_static_latency/LATENCY_SUMMARY_TABLE.csv}\newline \texttt{f94d0a7a9e95} & hardware-specific latency & \textit{Allowed:} Measured on the frozen RTX 4060 FP32 sequential workload. \newline \textit{Prohibited:} Portable latency, low cost, or universal serving efficiency. \\
\bottomrule
\end{tabular}
\end{table*}

\begin{table*}[htbp]
\centering
\caption{Claim-to-artifact trace and wording guardrails (part 2 of 2). Hashes are the first 12 hexadecimal characters of SHA-256 over the named frozen artifact.}
\label{tab:claim-trace-2}
\scriptsize
\begin{tabular}{>{\raggedright\arraybackslash}p{0.04\linewidth}>{\raggedright\arraybackslash}p{0.17\linewidth}>{\raggedright\arraybackslash}p{0.22\linewidth}>{\raggedright\arraybackslash}p{0.11\linewidth}>{\raggedright\arraybackslash}p{0.25\linewidth}}
\toprule
ID & Headline result & Authoritative artifact / SHA-256$_{12}$ & Evaluation type & Allowed wording; prohibited wording \\
\midrule
C8 & Chronos-2 gains -3.048\% versus Global Fixed and fails qualification. & \path{experiments/06_5_teacher_qualification/outputs/CHRONOS2_QUALIFICATION_SUMMARY.json}\newline \texttt{01e8084c16ae} & zero-shot internal screen & \textit{Allowed:} This locked revision does not qualify for this protocol. \newline \textit{Prohibited:} Foundation models are generally inferior for EV data. \\
C9 & Aligned KD gains 0.119\% versus identical GT-only; CI crosses zero. & \path{experiments/07_residual_distillation/outputs/D1_SUMMARY.json}\newline \texttt{c38163157c89} & internal real GT; gated & \textit{Allowed:} The primary distillation comparison fails; distill\_v1 stops. \newline \textit{Prohibited:} Distillation improves the student or merits D2 latency. \\
C10 & Aligned KD gains 0.259\% versus shuffled teacher. & \path{experiments/07_residual_distillation/outputs/D1_SUMMARY.json}\newline \texttt{c38163157c89} & negative control & \textit{Allowed:} Evidence of a small sample-aligned teacher signal. \newline \textit{Prohibited:} Engineering value or a passed primary KD claim. \\
C11 & Paris teacher gains 0.847\%, wins 4/4 folds and 16/16 cells, but fails 1\%. & \path{experiments/09_distill_v2/outputs/teacher_qualification/aggregate/PARIS_TEACHER_QUALIFICATION_SUMMARY.json}\newline \texttt{5b7c1ddc353c} & development real GT & \textit{Allowed:} Consistent development gain below the frozen qualification effect. \newline \textit{Prohibited:} Round to 1\%, imply formal/protected performance, or call it official TimeXer. \\
C12 & Three identity shuffles were blocked at 69/72 before metric unlock. & \path{experiments/07_residual_distillation/outputs/FORMAL_MATRIX_AUDIT_V1_0_FAIL_IDENTITY_SHUFFLE.json}\newline \texttt{584313534390} & artifact integrity & \textit{Allowed:} The lock forced a non-selective rerun of all 24 controls. \newline \textit{Prohibited:} Use the invalid controls as performance evidence. \\
C13 & Paris prediction audit passed 32/32 before development-metric unlock. & \path{experiments/09_distill_v2/outputs/teacher_qualification/PARIS_TEACHER_PREDICTION_AUDIT.json}\newline \texttt{ee882d8e907b} & artifact integrity & \textit{Allowed:} Identity checks authorized development-metric aggregation. \newline \textit{Prohibited:} The audit proves model correctness or improved RMSE. \\
C14 & Protected analytical-access counter is 0; vault materialization is true. & \path{experiments/08_paris_new_evidence/PROTECTED_ACCESS_GUARD.json}\newline \texttt{a7780e39f4ef} & procedural data boundary & \textit{Allowed:} No protected label was used analytically after the boundary lock. \newline \textit{Prohibited:} Targets never existed, were never materialized, or are cryptographically inaccessible. \\
\bottomrule
\end{tabular}
\end{table*}


\FloatBarrier

\subsection{Latency across the frozen batch grid}
\label{sec:appendix-latency-grid}
\label{app:latency-all}

The main text uses batch-one end-to-end P50 as its primary engineering endpoint.  Table~\ref{tab:latency-all-batches} lists the complete batch 1/8/32 grid from the frozen latency summary.  The script \path{experiments/10_audit_paper/scripts/generate_revision_appendix.py} checks all 18 method--batch rows and formats the stored values without rerunning the benchmark.

The hardware and software settings were Windows 10 build 26200, 16 physical/32 logical AMD CPU cores, approximately 16.8 GB RAM, an RTX 4060 Laptop GPU with driver 595.71, Python 3.10.8, Torch 2.5.1+cu121, CUDA 12.1, cuDNN 9.1, FP32, and single-threaded Torch/OMP/MKL host settings.  GPU clocks were observed per process rather than pinned.  Model loading and checkpoint I/O were excluded.  Cold-start, quantized, compiled, parallel-expert, and serverless settings were not tested.

\begin{table*}[htbp]
\centering
\caption{Frozen end-to-end synchronous latency at batch sizes 1, 8, and 32. P50 and P95 are milliseconds per request; throughput is forecast origins per second at P50. Every value is formatted directly from \texttt{LATENCY\_SUMMARY\_TABLE.csv}; no benchmark was rerun.}
\label{tab:latency-all-batches}
\footnotesize
\begin{tabular}{lrrrr}
\toprule
Method & Batch origins & E2E P50 (ms) & E2E P95 (ms) & Origins s$^{-1}$ (P50) \\
\midrule
DLinear & 1 & 1.166 & 1.763 & 857.9 \\
DLinear & 8 & 2.118 & 2.874 & 3776.9 \\
DLinear & 32 & 5.382 & 7.728 & 5945.7 \\
\addlinespace[2pt]
NLinear & 1 & 0.674 & 1.056 & 1483.7 \\
NLinear & 8 & 1.554 & 2.101 & 5148.0 \\
NLinear & 32 & 3.896 & 5.373 & 8213.7 \\
\addlinespace[2pt]
GBDT-Forecaster & 1 & 6.242 & 7.195 & 160.2 \\
GBDT-Forecaster & 8 & 40.188 & 50.261 & 199.1 \\
GBDT-Forecaster & 32 & 157.087 & 186.205 & 203.7 \\
\addlinespace[2pt]
CAPER & 1 & 7.934 & 10.798 & 126.0 \\
CAPER & 8 & 10.740 & 14.253 & 744.9 \\
CAPER & 32 & 19.753 & 24.266 & 1620.0 \\
\addlinespace[2pt]
TimeXer & 1 & 14.067 & 15.952 & 71.1 \\
TimeXer & 8 & 105.780 & 113.816 & 75.6 \\
TimeXer & 32 & 482.165 & 519.827 & 66.4 \\
\addlinespace[2pt]
Global Fixed CAPER+TimeXer & 1 & 20.940 & 33.123 & 47.8 \\
Global Fixed CAPER+TimeXer & 8 & 125.350 & 136.947 & 63.8 \\
Global Fixed CAPER+TimeXer & 32 & 479.634 & 510.574 & 66.7 \\
\bottomrule
\end{tabular}
\begin{minipage}{0.94\linewidth}\vspace{3pt}\footnotesize
Workload: one request produces all four horizons for 275 UrbanEV zones. Measurements used FP32, one RTX 4060 Laptop GPU, sequential expert execution, five processes, 50 warm-ups, and 1,000 timed iterations per process and method--batch cell. Throughput is not monotone across model families because the frozen serving graphs batch differently; these values should not be assumed to hold in other deployment settings.
\end{minipage}
\end{table*}


\FloatBarrier

\subsection{Protocol and environment artifact index}
\label{sec:appendix-artifact-index}

The following index lists existing controls.  Listing a file does not make missing environment information complete.  The latency and Chronos stages have explicit environment records.  Other stages retain protocols, source/configuration fingerprints, and run receipts but do not provide one cross-stage environment hash.  Exact environment equivalence is therefore unknown/not claimed where that record is absent.

\begin{table*}[htbp]
\centering
\caption{Protocol and environment artifact index (part 1 of 2).}
\label{tab:artifact-index}
\scriptsize
\begin{tabular}{>{\raggedright\arraybackslash}p{0.40\linewidth}>{\raggedright\arraybackslash}p{0.28\linewidth}>{\raggedright\arraybackslash}p{0.23\linewidth}}
\toprule
Artifact & Function & Scope limit \\
\midrule
\path{experiments/00_registry_reclassification/scripts/rebuild_status_registry.py} & Deterministic S1--S6 implementation and priority order & Single programmed classifier; no inter-rater validation \\
\path{experiments/00_registry_reclassification/outputs/CONFIG_STATUS.csv} & Row-level run IDs, fingerprints, grades, reasons, and rerun policy & One frozen UrbanEV inventory \\
\path{experiments/02_lightweight_baselines/configs/BASELINE_PROTOCOL.json} & UrbanEV target, folds, horizons, seed, and candidate protocol & Candidate-set internal evaluation only \\
\path{refine-logs/ROUTER_PROTOCOL_V1_0.md} and \path{experiments/05_router_v1/configs/ROUTER_V1.json} & Router feature/training configuration and frozen Gate 4 & One seed; proxy and real-GT roles remain separate \\
\path{experiments/06_static_latency/LATENCY_PROTOCOL_V1_0.md} & Request shape, synchronization, warm-up, repetitions, and primary endpoint & Warm synchronous serving, not cold start \\
\path{experiments/06_static_latency/hardware_manifest.json} and \path{experiments/06_static_latency/software_manifest.json} & Device, driver, platform, runtime, libraries, precision, host threads & One laptop system; clocks observed rather than pinned \\
\path{experiments/06_static_latency/LATENCY_SUMMARY_TABLE.csv} & Batch 1/8/32 model-only and end-to-end distributions, throughput, memory & Hardware/workload-specific \\
\path{experiments/06_5_teacher_qualification/TEACHER_QUALIFICATION_PROTOCOL_V1_0.md} & Chronos-2 qualification endpoint and gate & Exact project protocol only \\
\path{experiments/06_5_teacher_qualification/CHRONOS2_MODEL_LOCK.json} & Resolved model commit, file hashes, snapshot bytes, and environment & One revision and zero-shot profile \\
\path{experiments/07_residual_distillation/DISTILL_PROTOCOL_V1_0.md} & Teacher, student, loss, controls, milestones, stop rules & Distill v1 only; later search prohibited after failure \\
\path{experiments/07_residual_distillation/configs/FORMAL_STUDENT_CONFIG_V1.json} & 338-parameter student input/architecture/training identity & One architecture and seed \\
\path{experiments/07_residual_distillation/SHUFFLE_CONTROL_SPEC.json} & 24-hour block-shuffle negative-control construction & Control test, not a second primary endpoint \\
\bottomrule
\end{tabular}
\end{table*}

\FloatBarrier

\begin{table*}[htbp]
\centering
\caption{Protocol and environment artifact index (part 2 of 2).}
\label{tab:artifact-index-2}
\scriptsize
\begin{tabular}{>{\raggedright\arraybackslash}p{0.40\linewidth}>{\raggedright\arraybackslash}p{0.28\linewidth}>{\raggedright\arraybackslash}p{0.23\linewidth}}
\toprule
Artifact & Function & Scope limit \\
\midrule
\path{experiments/07_residual_distillation/outputs/FORMAL_MATRIX_AUDIT_V1_0_FAIL_IDENTITY_SHUFFLE.json} and \path{experiments/07_residual_distillation/outputs/FORMAL_MATRIX_AUDIT.json} & Pre-unlock failure and final 72/72 audit & Supports control integrity, not accuracy \\
\path{experiments/07_residual_distillation/outputs/D1_SUMMARY.json} and \path{experiments/07_residual_distillation/outputs/D1_DECISION.md} & Primary/negative-control effects and terminal gate decision & No D2 latency, multi-seed, or protected evidence \\
\path{experiments/08_paris_new_evidence/DATA_BOUNDARY_LOCK.json} and \path{experiments/08_paris_new_evidence/PROTECTED_ACCESS_GUARD.json} & Role dates, permitted metadata, storage state, access counter, unlock rule & Procedural guard; not hardened authorization \\
\path{experiments/08_paris_new_evidence/PARIS_EVIDENCE_VAULT_REPORT.md} & Vault identities, hashes, negative main-project scan, threat limits & Metadata/hash audit only; no content-level protected audit \\
\path{experiments/09_distill_v2/PARIS_TEACHER_QUALIFICATION_PROTOCOL_V1_0.md} & Development folds, target, missingness, time rule, teacher fit, gate & Paris development only \\
\path{experiments/09_distill_v2/PARIS_TEACHER_MODEL_CONFIG.json} & Local CAPER/TimeXer profiles, training, queue, bootstrap & Local adaptations, not exact upstream implementations \\
\path{experiments/09_distill_v2/TIMEXER_LOCAL_ADAPTATION_PROVENANCE.json} and \path{experiments/09_distill_v2/PARIS_CAPER_ADAPTATION_V1_0.md} & Source hash and explicit local adaptation differences & Upstream TimeXer commit undisclosed; Paris CAPER newly trained \\
\path{experiments/09_distill_v2/PARIS_TEACHER_DEPLOY_APPROVAL_V1_1.md} & Automated differential audit closure after the earlier block & Automated artifact audit, not blinded human replication \\
\path{experiments/09_distill_v2/outputs/teacher_qualification/PARIS_TEACHER_PREDICTION_AUDIT.json} and \path{experiments/09_distill_v2/outputs/teacher_qualification/PARIS_TEACHER_METRICS_UNLOCK_RECEIPT.json} & 32/32 closure, identities, role-specific metric unlock, access counts & Development-metric authorization only \\
\path{experiments/09_distill_v2/outputs/teacher_qualification/aggregate/PARIS_TEACHER_QUALIFICATION_SUMMARY.json} & Point effect, fold/cell results, synchronized bootstrap, gate vector & Current hypothesis failed; no student/formal/protected result \\
\path{experiments/09_distill_v2/PARIS_TEACHER_RESULT_REVIEW.md} & Separately instantiated result-to-claim check and provenance warnings & Warnings limit stronger process claims; they do not reverse the gate \\
\path{experiments/10_audit_paper/scripts/generate_revision_appendix.py} & Mechanical hash/latency table generation & No model, metric, resampling, or vault operation \\
\bottomrule
\end{tabular}
\end{table*}

\FloatBarrier

\subsection{Author and reviewer application checklist}
\label{sec:appendix-checklist}

Table~\ref{tab:audit-checklist} summarizes checks that authors and reviewers can apply.  Missing validation must be stated as missing.  The 1\% rule is a local convention without an economic utility model.  The S1--S6 labels come from one deterministic classifier and have no inter-rater evidence.  The model and latency conclusions are limited to their recorded seeds, panels, adaptations, and serving setup.

\begin{table*}[htbp]
\centering
\caption{Portable author and reviewer checklist for auditable forecasting claims.}
\label{tab:audit-checklist}
\scriptsize
\begin{tabular}{>{\raggedright\arraybackslash}p{0.035\linewidth}>{\raggedright\arraybackslash}p{0.27\linewidth}>{\raggedright\arraybackslash}p{0.30\linewidth}>{\raggedright\arraybackslash}p{0.29\linewidth}}
\toprule
ID & Check & Author must expose & Reviewer test \\
\midrule
1 & Target and role semantics & numerator/denominator, spatial unit, mask, sampling clock, and development/formal/protected role & Could two datasets or target constructions be mistaken for the same RMSE scale? \\
2 & Configuration identity & run ID, protocol version, fingerprint, source mapping, seed, fold/horizon, and execution state & Could a completed run still be underidentified or map to several disclosed cells? \\
3 & Information timing & forecast origin, lookback, covariates, fill/scaler fit interval, and prequential weight/tuning source & Does any current/future target influence a parameter evaluated in the same fold? \\
4 & Artifact closure & expected matrix cardinality, terminal attempts, prediction/target/order/time/spatial/mask identities & Are aggregates withheld until every required cell is closed and verified? \\
5 & Metric lock & metric implementation identity, lock state, unlock conditions, authority, receipt, and access counter & Can any result be read or selected before the audit passes? \\
6 & Comparator and endpoint & primary reference, macro/micro weighting, endpoint, sign convention, and secondary diagnostics & Was the reference or endpoint changed after seeing results? \\
7 & Uncertainty and gate & resampling unit, partial-block rule, replicates/seed, minimum effect, consistency checks, stop rule & Is a positive interval against zero being substituted for a failed effect-size requirement? \\
8 & Controls and repair & negative-control construction, invalidation scope, before/after hashes, and whether predictions changed & Were only visibly bad cells repaired, or was the formal control family rerun non-selectively? \\
9 & Model identity & upstream/local source, checkpoint/revision, local adaptations, input contract, and every unknown field & Does a local adaptation masquerade as an official or exact upstream implementation? \\
10 & Engineering measurement & device, software, precision, batch/request shape, synchronization, warm-up, repetitions, P50/P95, throughput & Is a hardware-specific measurement being promoted to portable efficiency? \\
11 & Future/protected evidence & storage versus analytical access, allowed metadata, threat-model limit, logs, and unlock rule & Does the boundary prevent analysis, merely record storage, or claim stronger security than implemented? \\
12 & Claim lineage & claim ID, authoritative artifact, SHA-256, evaluation type, allowed/prohibited wording, scope limit & Can every headline number be traced, and are diagnostics/subchecks kept at their original evidence type? \\
13 & Negative result & terminal decision, disallowed same-interval actions, and requirements for a new version/evidence boundary & Was a failed hypothesis silently continued through tuning, threshold changes, or a new structure? \\
14 & Missing validation & limitations must name absent multi-seed, multi-hardware, external-team, utility, or rater evidence & Is any absent analysis implied by rhetoric such as robust, validated, secure, or general? \\
\bottomrule
\end{tabular}
\end{table*}


\FloatBarrier
